\section{Decision Quality}

The agent is expected to achieve higher decision quality than the selected baseline policies in ambiguous incident scenarios (hypothesis H1). Specifically, it is expected to resolve a larger share of incidents correctly than the initial-symptom, scripted and always-escalate policies at equal or lower investigation cost, and to approach the confidence-threshold policy without being given the likelihood model. Numerical targets will be fixed after the pilot runs, and the results will be reported whether or not the hypothesis holds.

\section{Rule Compliance}

With the safety mechanism in place, the agent is expected to achieve a higher rate of rule-compliant correct recommendations than the same agent without it (hypothesis H2). No prohibited action will be executed when the mechanism is active, so the comparison of interest is the compliant success rate against the unconstrained and penalty-only agents. The proposed-violation rate is expected to fall over training, which would show that the agent learns to act within the rules and does not rely on being stopped.

\section{Robustness to Ambiguity}

As the level of ambiguity rises, the agent is expected to request more diagnostic information and to escalate more incidents, and not to continue recommending fixes on weaker evidence (hypothesis H3). The agent trained across randomised ambiguity levels is expected to degrade less than the agent trained at a single level. The experiments will report what no reviewed study has reported: how the shares of investigating, recommending and escalating change across controlled levels of missing, unreliable and overlapping evidence.

\section{Prototype}

The research is expected to produce a working simulator of IT incident triage built from the Directorate's scenarios, the rule-based safety mechanism, six baseline policies, trained agents, and scripts that reproduce every experiment. A lightweight dashboard will replay an incident step by step, showing the evidence revealed, the agent's estimated action values, any blocked proposal and the final decision, beside the comparison results.

\textbf{Broader impact.} Putting this together, the expected results address the problem this research is concerned with: that automated triage tools predict an answer for every incident and do not know when to look further, when to act within the rules, and when to hand the incident to a person. An agent that makes these three decisions under explicit, readable rules, evaluated on all four groups of measures, would give IT support teams such as the Directorate evidence on how far such automation can be trusted before it is considered for use.
